\Titular*%
{divulgacion}%
{Aplicacións da física no campo da medicina}%
{Daniel Bostán Boxenean}%
{Técnicas de imaxe por Resonancia Magnética Nuclear.}%

\begin{multicols}{2}

\subsection*{Introdución}

Acode ao Servizo de Neuroloxía un varón de 56 anos, fumador dun paquete diario
durante 34 anos. Diagnosticado de carcinoma microcítico de lobo pulmonar medio
(dereito) hai 15 meses, sendo positivo para PD-L1 e CD-56, parece responder ao
tratamento tras oito sesións e catro meses de carboplatino e etopósido a doses
elevadas segundo protocolo NCCN 2024. Na consulta refire hemiparesia
contralateral e problemas na fala, polo que o relacionamos a afectación de
cápsula interna e córtex. Exploración física: PA: 126/87 mmHg, FC: 85 lpm,
SatO2: 98\%, auscultación pulmonar e cardíaca normais. Semella consciente,
afebril. Porén, sospeitamos lesión do Sistema Nervioso Central, dado que
presenta hiperreflexia e signo de Babinski positivo, así que derivámolo a
radioloxía ante sospeita de metástase (fig. \ref{dbb:1}).

\begin{figure}[H]
    \centering
    \includegraphics[width=0.65\linewidth]{revistas/003/imaxes/dbb1.png}
    \caption{
        RMN en corte axial T1. Obsérvase unha masa tumoral
        no lobo temporal dereito, con afectación talamocortical.
    }
    \label{dbb:1}
\end{figure}

Pero de que maneira Philips ou Siemens\footnote{Xunto con General Electrics e
Toshiba, as principais comercializadoras de equipos NMR hoxe en día.} conseguen
obter tal detalle cun \textit{doughnut} de metal xigante?

\subsection*{Historia da Resonancia Magnética Nuclear}

En 1938, Rabi bautiza o fenómeno da Resonancia Magnética Nuclear (NMR),
consistente en que algúns núcleos teñen a capacidade de absorber e remitir
enerxía en forma de radiofrecuencia (RF) se están sometidas a un campo
magnético $\vec{B}_0$. Non foi ata 1973 que Paul Lauterbur conseguiu un método
para codificar espacialmente os sinais da NMR empregando gradientes lineais de
campo. A diferenza doutras técnicas de imaxe convencionais, a NMR non emprega
radiación ionizante --as ecografías tampouco, pois interpretan sinais por
ultrasóns--, o que a fai unha alternativa inocua.

\subsection*{Concepto básico da NMR}

Na NMR diagnóstica, emprégase o átomo de hidróxeno pola sinxela razón de que os
animais temos abundante H, principalmente en forma de auga e graxa (85\% aprox.
da masa total). Ademais, a procura do protón permite distinguir os diferentes
tecidos que atopamos no corpo debido ás propiedades químicas que teñen.

O momento angular do núcleo magnetizado produce un momento magnético
($\vec{\mu}$), definido pola relación xiromagnética ($\gamma$) particular para
cada átomo (42,58 MHz/T no hidróxeno). Como é de esperar, un conxunto de
momentos magnéticos non terá aliñación ningunha ata que apliquemos un campo
magnético externo ($\vec{B}_0$), o que fai que os spins se ordenen co campo
magnético (fig. \ref{dbb:2} esquerda). É fundamental para entender a NMR
coñecer a frecuencia de Larmor $\omega_0$. (fig. \ref{dbb:2} dereita)

\begin{figure}[H]
    \centering
    \begin{minipage}{0.3\linewidth}
        \centering
        \includegraphics[scale=0.75]{revistas/003/imaxes/dbb3.png}
    \end{minipage}
    \begin{minipage}{0.6\linewidth}
        \centering
        \includegraphics[scale=0.93]{revistas/003/imaxes/dbb2.png}
    \end{minipage}
    \caption{
        Esquerda: Núcleo aliñado a un campo magnético. Dereita: Núcleo en fase
        de relaxación lonxitudinal $T1$. Ecuación de Larmor:
        $\vec{\omega}=\gamma\vec{B}$
    }
    \label{dbb:2}
\end{figure}

Xa entendemos por que nuclear, e tamén por que magnética, pero por que
resonancia? Os spins están aliñados --case-- perfectamente con $\vec{B}$, pero
ao ``atordalos'' cun pulso de radiofrecuencia (pulso RF) seguirán xirando en
sentido perpendicular ao campo magnético coa mesma frecuencia $\omega$ que o
pulso RF, e esa oscilación emite un sinal medible. Pasado un tempo (variable
segundo o tecido e o pulso RF), o spin volve a orientarse con $\vec{B}$
(denominado magnetización ou relaxación lonxitudinal). O movemento que realiza
durante a reorientación é a precesión (e nutación) arredor de $\vec{B}$ cun
ángulo $\theta$, que é de 90º con cada pulso RF, e decae coa relaxación
lonxitudinal, xa que lembremos que en todo momento o spin segue a rotar coa
mesma frecuencia. Veremos en profundidade a interpretación dos compoñentes
vectoriais e o que significan (fig. \ref{dbb:3}).

\begin{equation}
    S(t)=\gamma NB\sin{(\theta)}\cos{(\omega t)}
\end{equation}

\begin{figure}[H]
    \centering
    \includegraphics[width=0.38\linewidth]{revistas/003/imaxes/dbb4.png}
    \caption{
        En gris $\vec{B}$, en vermello $\vec{\mu}$, en verde a proxección de
        $\vec{\mu}$, que constitúe a intensidade do sinal ao introducilo nunha
        bobina.
    }
    \label{dbb:3}
\end{figure}

Lembremos algúns fundamentos como a lei de Lenz, entre outros, para ter ben en
conta os determinantes da amplitude e frecuencia da función. Debemos entón
diferenciar moi ben dous fenómenos:

\begin{itemize}
    \item[--] Relaxación $T2$: é o tempo que o núcleo rota pola
        influencia do pulso RF perpendicular ao campo magnético.
        Durante $T2$, $\theta$ é máximo e constante. O sinal
        obtido do conxunto de protóns denomínase magnetización transversal.
    \item[--] Relaxación $T1$: entre o instante en que remata
        $T2$ e se aliña correctamente con $\vec{B}$. En $T1$,
        $\theta$ vai decaendo. O sinal do conxunto é a
        magnetización lonxitudinal.
\end{itemize}

Cada tecido ten uns tempos de relaxación distintos, incluso pode ser que
$T1\neq T2$ no mesmo tecido, por exemplo en $T2$ veremos o líquido
cefalorraquídeo de cor branca, mentres que en $T1$ será negro (fig.
\ref{dbb:1}), polo que cada zona emitirá un sinal distinto, que se traduce
nunha imaxe de altísima precisión.

\subsection*{Relaxación $T2$}

Agora ben, non estamos tendo en conta a cantidade innumerable de átomos de $H$
que hai por unidade de volume; e aínda que comezan (tras o pulso RF) co mesmo
momento angular, pasado un tempo empezan a desfasarse e decaendo $\theta$, o
que dá como resultado unha curva de “\textit{decay}” na amplitude desas
oscilacións no eixo horizontal (fig. \ref{dbb:4}). Hai basicamente dous
parámetros que determinan o contraste entre os tecidos. Un deles é ese
\textit{decay}, onde cada tecido ten un $T2$ distinto (curva distinta). O outro
é $N$, é dicir, a densidade de núcleos xirando (a amplitude correspondente a
$t=0$). Estas características son intrínsecas do tipo de tecido que estamos a
estudar. Porén, o que o radiólogo pode modular é o \textit{Time of Echo} (TE)
que é o tempo que esperamos para recibir o sinal dende que se enviou o pulso
RF.

\begin{figure}[H]
    \centering
    \includegraphics[width=0.65\linewidth]{revistas/003/imaxes/dbb5.png}
    \caption{
        Funcións \textit{decay} en $T2$ para distintos tecidos. O sinal,
        loxicamente, é sinusoidal e decae, polo que a curva da imaxe obtémola
        de axustar a ecuación de forma que vemos a variación das amplitudes no
        tempo. Imaxe NMR de escala de grises invertida, recurso moi empregado
        para resaltar mellor estruturas como a duramáter.
    }
    \label{dbb:4}
\end{figure}

\subsection*{Relaxación $T1$}

A curva de $T1$ é distinta pero intuitiva. É, por así dicilo, o fenómeno
inverso ou consecuente de $T2$. Estando todos os núcleos xirando con ángulos
distintos, xa desfasados, pasado ese tempo ($T1$) a maioría dos núcleos están
remagnetizados con $\vec{B}$. Neste caso, obtemos unha gráfica característica
que, como xa veremos, complementa ben o sinal obtido na relaxación $T2$ (fig.
\ref{dbb:5}). Podemos tamén modular o tempo que esperamos entre estimulacións
repetidas, é dicir, cada canto enviamos un pulso RF. É o que chamamos
\textit{Time of Repetition} (TR)

\begin{figure}[H]
    \centering
    \includegraphics[width=0.7\linewidth]{revistas/003/imaxes/dbb6.png}
    \caption{
        Curvas de $T1$, coa barra gris indicando o TR. Pódese observar que un
        tecido con $T1$ longo (rosa) corresponde cunha zona con menos sinal
        para ese TR. Imaxe con cor invertida.
    }
    \label{dbb:5}
\end{figure}

\subsection*{Relacións das relaxacións}

O desafío está en escoller os TE e TR adecuados para cada situación, que
depende non só da “exposición” que queiramos (como se dunha fotografía se
tratase), senón das características do propio tecido. O rabdomiosarcoma é un
tumor mesenquimal\footnote{Mesénquima: tecido de orixe mesodérmico, sumado a
endodermo e ectodermo, que son as tres capas xerminais de calqueira tecido
embrionario con capacidade de diferenciación, neste caso, a tecido conectivo
(músculo, óso, graxa, cartílago…).} de orixe embrionaria que tende a
diferenciarse a músculo esquelético. Ante un estudo por NMR, o radiólogo ten
que ter en conta a exploración física e a historia clínica para orientar un
diagnóstico e así predicir non só o tecido que buscamos, senón o tipo de tumor
(que como ten orixe mesodérmica, será semellante ao músculo circundante), polo
que debe elixir as secuencias máis óptimas en función dos $T1$ e $T2$ da masa
que buscamos.

\begin{figure}[H]
    \centering
    \includegraphics[width=0.75\linewidth]{revistas/003/imaxes/dbb7.png}
    \caption{
        Relación TR/TE na recepción de sinal. Un TR curto
        inflúe na amplitude en $t=0$ para a fase $T2$ (liña laranxa).
        Cada nova barra laranxa é un pulso RF, o que ``reinicia'' un $T2$.
    }
\end{figure}

Como se pode comprobar na Figura 7, o TE depende directamente do TR, e con el o
sinal que esperamos recibir dos tecidos. Un tecido hipotético cun $T1$ máis
curto, cuxa curva estaría enriba da azul, terá unha amplitude maior na curva
laranxa en $t=$0. O pulso RF é como pulsar \textit{play} nun cronómetro, que dá
comezo a $T2$. Se poñemos un TR e TE moi curtos, a relaxación transversal
($T2$) non nos daría moita información, xa que estariamos constantemente
excitando eses núcleos, e sempre estariamos nunha zona da gráfica con pouca
dispersión das curvas (porque cada tecido ten unha curva distinta). É por iso
que cando unha secuencia establece un TE e TR curtos, dicimos que está
potenciada ou ponderada en $T1$\footnote{Se poñemos un TR longo e un TE curto
obteriamos o que se chama Proton Density Image.}, e viceversa. Non son
excluíntes, senón complementarias.

Por suposto, $T1+T2$ constitúe o tempo que tarda o núcleo en aliñarse
totalmente con $\vec{B}$ dende o pulso, e aínda que medimos ambos os tempos
como fenómenos distintos, a realidade é que son un contínuum, e, por tanto, o
vector de magnetización pódese descompoñer en dous vectores perpendiculares. A
este sistema chamámoslle \textit{rotating frame}, e vemos como ao longo do
tempo, a compoñente horizontal decae (corresponde coa magnetización
transversal, que ocorre durante $T2$) e o vertical medra. Recomendo indagar
nestes aspectos, que carecen de interese clínico e, por tanto, non teñen lugar
neste artigo.

\subsection*{Polarización}

Todos os camiños levan a Boltzmann. Esta ecuación determina que fracción dos
núcleos están orientados de forma paralela e cantos de forma antiparalela, que
con parámetros corporais resulta un número moi baixo. Nun conxunto de 99 spins,
50 arriba e 49 abaixo, a inmensa maioría cancelará os seus momentos magnéticos,
e desta forma un só spin é o que emite sinal real aproveitable. No entanto,
lembro que o sinal depende tamén das características intrínsecas de cada
tecido, e que estamos omitindo cuestións como a non-homoxeneidade do campo
magnético, o spin \textit{echo} e moitas transformadas de Fourier.

\begin{equation}
    S=\frac{N_{total}\gamma^3\hbar^2 B_0^2}{4kT}
\end{equation}

Existe unha técnica baseada en hiperpolarizar gases nobres como o Xe-129 ou
He-3, dado que emiten moito máis sinal. Desta forma conseguimos unha mellor
resolución en determinadas estruturas, útil, por exemplo, en estudos en
asmáticos. Outro contraste útil (intravenoso) é o gadolinio.

\subsection*{Do imán á pantalla}

Para rematar, sabemos que a radiografía tradicional é unha proxección, e que
unha tomografía son radiografías en serie que ``cortan'' fisicamente nun plano.
A información recibida dunha NMR non é en planos; e, por tanto, non falamos de
píxeles, senón de vóxeles; aínda que o plano de referencia é o campo magnético,
que sempre é o mesmo\footnote{NUNCA se apaga un equipo NMR. Se pasades ao lado
do CACTUS escoitaredes un son periódico, “hidráulico”. É un sistema de
refrixeración de helio líquido, para manter a intensidade día e noite (sae máis
barato que apagalo).} (SEMPRE a 2-7 T, incluso cando non está funcionando a
máquina). Por defecto, obtemos imaxes no plano axial, pero as bobinas de
gradiente (as responsables dese insoportable e característico ruído)
``desprazan'' $\vec{B}$ de forma que, sen mover o equipo nin o paciente,
podemos obter imaxes en tempo real en calquera plano.

A NMR cambiou por completo o enfoque diagnóstico na medicina, e conclúe unha
proba máis da importancia na colaboración entre disciplinas científicas para a
saúde.

\nocite{harrison_2020}
\nocite{cunn_74}

\end{multicols}
